\scanpage{63}\mspage{62}
Donc

\uline{Corollaire} Pour un atelier ensembliste

\noindent\begin{tabular}{@{}|p{.92\linewidth}@{}}
(satisfaisant \deleted{\unsure{At 8|At}} aussi At 9 et Atens 2)
Si $F,G$ sont deux figures, pour que
$F\parallel G$, il \unsure{f|f.} et S que $\forall x\in\operatorname{omb}X,y\in\operatorname{omb}Y$,
on ait $x\parallel y$. \deleted{\unsure{ou encore}}
\end{tabular}

Dans le cas des ateliers \unsure{quelconque|quelconques}, on
trouve d'autre part.

\noindent\begin{tabular}{@{}|p{.92\linewidth}@{}}
\uline{Corollaire} Validité de $(1.58)$, si $\unsure{\mathcal F|F}$
satisfait At 9
\end{tabular}

D'autre part, je voudrais en présenter
un axiome facultatif

% slanted
\margin{\uline{NB} dans $\mathcal F$, le $\operatorname{Inf}$ \deleted{\unsure{de non|d'non}} vide existe tjrs, \unsure{$=$} \deleted{\unsure{pour}} $F\cap G$ sur $F,G$ \unsure{compatibles|compatible}, il correspond à l'intersection des parties \deleted{de $M$ de} fermées de $M$ de la forme $\widetilde F$. L'énoncé ci-contre s'étend \unsure{automatiquement} au cas de deux figures quelconques\ldots}
\uline{At \deleted{\unsure{ens fig|ensfiltr}}} \unsure{9$'$|91} Soient $X,Y\in M$, \deleted{$\unsure{F|E}$} et
soit \deleted{\unsure{au sens|on aura}} $F=X\cap Y$ \inserted{(donc $\widetilde F$ est} l'ens des strates communes).
\deleted{\unsure{qui est une partie fermée de} $M$,}
\inserted{Supposons que pour tout}
$x\in\operatorname{omb}(X)\unsure{\setminus|-}\operatorname{omb}(F)$ et tout $y\in\operatorname{omb}(Y)-\operatorname{omb}(F)$
on ait $x\unsure{\mathrel{\tikz[baseline=-.46ex,x=.68em,y=.68em]{\draw[line width=.075em] (.16,.64)--(.18,-.23) (.36,.67)--(.39,-.24) (.15,.23)--(.55,.22);}}|\parallel|\nparallel}y$ (donc $x\ne y$) $[$ce \unsure{que|qui}
\unsure{explique|implique} encore que \deleted{\unsure{on ait}}
\[
\operatorname{omb}(X)\cap\operatorname{omb}(Y)=\operatorname{omb}(F)\ \bigr]
\]
\diagram{\begin{tikzcd}[ampersand replacement=\&,column sep=small,row sep=small]
\operatorname{om}(X)\arrow[dr] \& \& \operatorname{omb}(Y)\arrow[dl,bend right=10]\\
\& \operatorname{omb}(F) \&
\end{tikzcd}}
Alors $X \mathrel{\tikz[baseline=-.45ex,x=1em,y=1em]{\draw[line width=.055em] (.52,.27)--(.05,.055)--(.48,-.11); \draw[line width=.055em] (.14,.21)--(.46,.04)--(0,-.17);}} Y$
\margin{\tikz[baseline=(note.base)]{\node[inner sep=2pt,outer sep=0pt,text width=\dimexpr\marginparwidth-4pt\relax,align=left] (note) {\uline{NB} la réciproque \deleted{\unsure{sera|être}}\unsure{reste|sera} vraie aussi, par axiome de disjonction plus bas p.72}; \draw[rounded corners=4pt] (note.north west)--(note.south west)--(note.south east);}}

Cet axiome est \unsure{``|''}\deleted{\unsure{in}}\unsure{intéressant}, en ce
qu'il exprime \deleted{\unsure{les|la}} compatibilité \inserted{$\mathrel{\tikz[baseline=-.45ex,x=1em,y=1em]{\draw[line width=.055em] (.52,.27)--(.05,.055)--(.48,-.11); \draw[line width=.055em] (.14,.21)--(.46,.04)--(0,-.17);}}$} des multistrates
génér\unsure{ales|iques} (qui \inserted{\unsure{pratiquement}} est un des \deleted{\unsure{grands points d'interrogation}}
